\scanpage{617}\mspage{5}
% archive sheet number: 616

égalité, faute de pouvoir montrer que $\Phi''\unsure{\natural|{}^{\natural}|\sharp}$
est stable, non seulement par facteurs directs, mais aussi par cochgt.
de base et par compositions transfinies.

Cependant, rappelons

\uline{\deleted{\unsure{Proposition}}\inserted{Théorème} 1.1}.
Supposons $M$ stable par petites $\varinjlim$, et que ses éléments
soient accessibles. Supposons de plus que $\Phi$ soit petit, plus
généralement qu'il existe une partie $\Phi_0\subset\Phi$ petite telle
que $\Phi\subset\Phi_0^{\sim}$, \deleted{Alors}\inserted{Soit}
\deleted{\unsure{Soit $\Phi$}}
\[
 \Psi\overset{\mathrm{def}}{=}\Phi_*=(\Phi_0)_* .\tag*{(1.1.1)}
\]
\margin{\begin{tabular}{@{}|l@{}}
\uline{NB} On a donc\\
$\Phi \mathrel{\tikz[baseline=-.45ex,x=1em,y=1ex]{\draw (0,0)--(2.15,0) (0.32,-0.5)--(0,0) (1.83,0.5)--(2.15,0) (1.83,-0.5)--(2.15,0);}} \Psi$\\
a fortiori\\
$\Phi_0 \mathrel{\tikz[baseline=-.45ex,x=1em,y=1ex]{\draw (0,0)--(2.15,0) (0.32,-0.5)--(0,0) (1.83,0.5)--(2.15,0) (1.83,-0.5)--(2.15,0);}} \Psi$\end{tabular}}

Sous ces conditions

$\Phi''_0$ désigne le Q-saturé \inserted{à g.} transfini
\deleted{\unsure{de}} de $\Phi_0$, alors on a
$\Phi''_0 \mathrel{\tikz[baseline=-.45ex,x=1em,y=1ex]{\draw (0,0)--(2.6,0) (0.32,-0.5)--(0,0) (2.28,0.5)--(2.6,0) (2.28,-0.5)--(2.6,0) (1.3,-.52)--(1.3,.52);}} \Psi$, i.e. \unsure{toute|\deleted{tout}} flèche $f$ de $M$ se
factorise en $f=pi$, avec $i\in\Phi''_0$, $p\in\Psi$.
% slanted
\margin{Il y a $=$ factorisation \uline{fonctorielle} en $f$,
\deleted{\unsure{pro}} propriété qui mérite une notation, \unsure{spéciale|cf. p. suiv.|cf. p. 64} $\Phi''_0 \mathrel{\tikz[baseline=-.45ex,x=1em,y=1ex]{\draw (0,0)--(2.2,0) (0.32,-0.5)--(0,0) (1.88,0.5)--(2.2,0) (1.88,-0.5)--(2.2,0) (1,-.52)--(1,.52) (1.17,-.52)--(1.17,.52);}} \Psi$}
A fortiori, si $\Phi''$ désigne le Q-saturé \deleted{\unsure{à}} \unsure{à g.|g.} transfini de $\Phi$, on a \deleted{$\unsure{\Phi}$}
\[
 \Phi'' \mathrel{\tikz[baseline=-.45ex,x=1em,y=1ex]{\draw (0,0)--(3,0) (0.32,-0.5)--(0,0) (2.68,0.5)--(3,0) (2.68,-0.5)--(3,0) (1.5,-.52)--(1.5,.52);}} \Psi .
\]

\deleted{\illegible{1}}\uline{Corollaire 1.2} : Sous les conditions
du th, on a
\[
 \Phi'''\ \substack{\text{(\deleted{\unsure{le|Q}}Q-saturé \inserted{à g.} total}\\
                      \text{de }\Phi\text{)}}
 =\underset{\underbrace{\substack{\text{ens des}\\\text{fact. directs}\\
             \text{d'éléments}\\\text{de }\Phi^{\unsure{\prime\prime|o}}}}}
          {\Phi''\natural}
 =\Phi^{\sim} .\tag*{(1.2.1)}
\]

\deleted{\unsure{P}}Plus généralement (en apparence seulement !)

Si $\Phi_1$ est une partie de $\mathrm{Fl}(M)$ qui est Q-saturée à
gauche transfinie, et telle que
\[
 \Phi_0\subset\Phi_1\subset\Phi^{\sim}=\Phi_0^{\sim},
\]
